\section{Preliminaries}
\label{sec:preliminaries}

\textcolor{blue}{
We begin with a brief overview of the contact-dynamics and contact-planning formulations used throughout this work.
Section~\ref{preliminaries:contact_dynamics} introduces time-stepping contact dynamics and the associated complementarity formulation.
Section~\ref{preliminaries:contact_planning} introduces model-based contact planning and reviews CIMPC as a representative online formulation.
Section~\ref{preliminaries:complementarity_free} reviews the complementarity-free contact model that motivates the lightweight contact-model approximation used later in CSO.
}


\subsection{Contact Dynamics}
\label{preliminaries:contact_dynamics}

Contact dynamics provides a physical description of interactions among the robot, object, and environment.
Here, we consider the commonly used Anitescu model~\cite{anitescu2006optimization}, which is formulated as the following time-stepping equation:
\begin{subequations}
\label{equ.quasi_dyn}
\begin{align}
    \epsilon\boldsymbol{M}_o\boldsymbol{v}_o
    &=
    h\boldsymbol{\tau}_o
    +
    \sum_{i=1}^{n_c}
    \boldsymbol{J}_{o,i}^{\top}\boldsymbol{\beta}_i,
    \label{equ.quasi_dyn_1}
    \\
    h\boldsymbol{K}_r
    \left(
    h\boldsymbol{v}_r-\boldsymbol{u}
    \right)
    &=
    h\boldsymbol{\tau}_r
    +
    \sum_{i=1}^{n_c}
    \boldsymbol{J}_{r,i}^{\top}\boldsymbol{\beta}_i.
    \label{equ.quasi_dyn_2}
\end{align}
\end{subequations}

\textcolor{blue}{
Here, $\boldsymbol{\beta}_i$ denotes the contact impulse over the time interval $h$.
In the contact-force formulation below, we use $\boldsymbol{\lambda}_i$ to denote the corresponding contact force.
Assuming that the contact force is approximately constant within one time step, they are related by
$\boldsymbol{\beta}_i \approx h\boldsymbol{\lambda}_i$.
}
$\epsilon\boldsymbol{M}_o$, $\boldsymbol{v}_o$, $\boldsymbol{\tau}_o$, and $\boldsymbol{J}_{o,i}$ are the regularized mass matrix, velocity change, non-contact force, and contact Jacobian of the object, respectively.
$\boldsymbol{K}_r$, $\boldsymbol{v}_r$, $\boldsymbol{\tau}_r$, and $\boldsymbol{J}_{r,i}$ are the robot stiffness matrix, velocity, non-contact torque, and contact Jacobian, respectively.
\textcolor{blue}{
The control input is defined as $\boldsymbol{u}=\Delta\boldsymbol{q}$.
}

The above equations describe, from the perspective of momentum balance, how the robot influences the object's state through contact, as well as how the object interacts with the environment.
Here, we follow~\cite{pang2021convex,jin2024complementarity} and consider the robot to be under impedance control~\cite{hogan1984impedance}, modeling the robot--object interaction as a spring system.
\textcolor{blue}{
The contact constraints between contact forces and system motion are typically formulated as the following cone complementarity problem~\cite{anitescu2010iterative}:
}
\begin{equation}
    \mathcal{K}_i
    \ni
    \boldsymbol{\lambda}_i
    \perp
    \boldsymbol{J}_i\boldsymbol{v}
    +
    \frac{1}{h}
    \begin{bmatrix}
        \phi_i \\
        0 \\
        0
    \end{bmatrix}
    \in
    \mathcal{K}_i^{*},
    \qquad
    \forall i\in\{1,\ldots,n_c\},
    \label{equ.complementarity_constraint}
\end{equation}
where
    $\mathcal{K}_i
    =
    \left\{
    \boldsymbol{\lambda}_i\in\mathbb{R}^{3}
    \,\middle|\,
    \mu_i\lambda_i^{n}
    \geq
    \left\|\boldsymbol{\lambda}_i^{d}\right\|
    \right\}$
is the Coulomb friction cone, and $\mathcal{K}_i^{*}$ is the dual cone as
    $\mathcal{K}_i^{*}=\left\{
    \boldsymbol{v}
    \,\middle|\,
    \boldsymbol{J}_i^{n}\boldsymbol{v}
    +
    \frac{\phi_i}{h}
    \geq
    \mu_i
    \left\|
    \boldsymbol{J}_i^{d}\boldsymbol{v}
    \right\|
    \right\}$.

Equations~\eqref{equ.quasi_dyn} and~\eqref{equ.complementarity_constraint} are the Karush--Kuhn--Tucker (KKT)~\cite{avriel2003nonlinear} optimality conditions of the following primal optimization~\cite{anitescu2006optimization}:
\begin{subequations}
\label{equ.contact_model}
\begin{align}
    \min_{\boldsymbol{v}}\quad
    &
    \frac{1}{2}h^{2}
    \boldsymbol{v}^{\top}
    \boldsymbol{Q}
    \boldsymbol{v}
    -
    h
    \boldsymbol{v}^{\top}
    \boldsymbol{b}(\boldsymbol{u}),
    \label{equ.qs_model_obj}
    \\
    \mathrm{s.t.}\quad
    &
    \boldsymbol{J}_i\boldsymbol{v}
    +
    \frac{1}{h}
    \begin{bmatrix}
        \phi_i\\
        0\\
        0
    \end{bmatrix}
    \in
    \mathcal{K}_i^{*},
    \qquad
    i\in\{1,\ldots,n_c\}.
    \label{equ.contact_model_con}
\end{align}
\end{subequations}
where
$\boldsymbol{Q} \in
\mathbb{R}^{(n_o+n_{\text{DoF}})\times(n_o+n_{\text{DoF}})}$
and
$\boldsymbol{b}(\boldsymbol{u}) \in
\mathbb{R}^{\textcolor{blue}{n_o+n_{\text{DoF}}}}$
are
\begin{equation}
    \boldsymbol{Q}
    :=
    \begin{bmatrix}
        \epsilon\boldsymbol{M}_o/h^{2} & \boldsymbol{0}\\
        \boldsymbol{0} & \boldsymbol{K}_r
    \end{bmatrix},
    \qquad
    \boldsymbol{b}(\boldsymbol{u})
    :=
    \begin{bmatrix}
        \boldsymbol{\tau}_o\\
        \boldsymbol{K}_r\boldsymbol{u}+\boldsymbol{\tau}_r
    \end{bmatrix}.
    \label{equ.q_matrix}
\end{equation}
Here, $n_o$ and $n_{\mathrm{DoF}}$ are the dimensions of the object and robot states, respectively, and $n_c$ is the number of contacts.

Problem~\eqref{equ.contact_model} is a second-order cone program (SOCP) and is often approximated by a QP through a polyhedral approximation of the friction cone:
\begin{equation}
    \left(
    \boldsymbol{J}_i^{n}
    -
    \mu_i\boldsymbol{J}_{i,j}^{d}
    \right)
    \boldsymbol{v}
    +
    \frac{\phi_i}{h}
    \geq 0,
    \quad
    i\in\{1,\ldots,n_c\},
    \quad
    j\in\{1,\ldots,n_d\},
\end{equation}
where $n_d$ is the number of friction-cone facets.
Solving~\eqref{equ.contact_model} corresponds to one step of physically consistent contact simulation.
The resulting velocity $\boldsymbol{v}$ can then be used to integrate the current state as
$\boldsymbol{q}_{t+1}
=
\boldsymbol{q}_{t}
\oplus
h\boldsymbol{v}$.
Problem~\eqref{equ.contact_model} and its variants are widely used in physics simulators.


\subsection{Model-based Contact Planning}
\label{preliminaries:contact_planning}

Contact planning refers to generating trajectories for approaching, making contact, and breaking contact with the object to manipulate it to a target state.
Model-based contact planning methods integrate contact dynamics and solve trajectory optimization problems.
As a representative contact planning method, CIMPC is formulated as the following optimization problem:
\begin{equation}
\begin{aligned}
    \min_{\boldsymbol{u}_{0:N-1}}\quad
    &
    \sum_{k=0}^{N-1}
    \ell(\boldsymbol{x}_k,\boldsymbol{u}_k)
    +
    V_f(\boldsymbol{x}_N)
    \\
    \mathrm{s.t.}\quad
    &
    \boldsymbol{x}_{k+1}
    =
    f(\boldsymbol{x}_k,\boldsymbol{u}_k),
    \quad
    k=0,\ldots,N-1,
    \\
    &
    \boldsymbol{x}_k\in\mathcal{X},
    \quad
    k=0,\ldots,N,
    \\
    &
    \boldsymbol{u}_k\in\mathcal{U},
    \quad
    k=0,\ldots,N-1.
\end{aligned}
\label{eq:CIMPC_opt}
\end{equation}
in which $\boldsymbol{x}_k$ is the system state, including the object pose $q_k$ and robot configuration,  and $f$ denotes the contact dynamics.
The stage cost $\ell$ can take various forms and is typically used to guide the process of approaching the contact location.
The terminal cost $V_f$ is usually defined as the error between the object pose and the desired pose:
\begin{equation}
\label{equ:obj_pose_cost}
\ell_{\mathrm{pose}}
=
w_{\mathrm{pos}}
\left\|
\boldsymbol{p}_{\mathrm{obj}}
-
\boldsymbol{p}_{\mathrm{goal}}
\right\|_2^2
+
w_{\mathrm{quat}}
\left(
1-
\left(
\boldsymbol{\psi}_{\mathrm{obj}}^{\top}
\boldsymbol{\psi}_{\mathrm{goal}}
\right)^2
\right).
\end{equation}

\textcolor{blue}{
The hybrid nature of contact dynamics renders~\eqref{eq:CIMPC_opt} a piecewise-smooth nonlinear complementarity problem (NCP).
Directly using the full contact model in~\eqref{equ.contact_model} within receding-horizon trajectory optimization remains computationally demanding, motivating the efficient contact-model approximations introduced below.
}


\subsection{Complementarity-free Model}
\label{preliminaries:complementarity_free}

\textcolor{blue}{
We next introduce the complementarity-free contact model proposed in~\cite{jin2024complementarity}, which motivates the surrogate contact-model approximation used later in CSO.
}
Problem~\eqref{equ.contact_model} has the dual form:
\begin{multline}
\label{equ.dual_prob_relax}
    \max_{\boldsymbol{\beta}\geq\boldsymbol{0}}
    \quad
    -\frac{1}{2h^{2}}
    \boldsymbol{\beta}^{\top}
    \left(
    \boldsymbol{\widetilde{J}}
    \boldsymbol{Q}^{-1}
    \boldsymbol{\widetilde{J}}^{\top}
    +
    \boldsymbol{R}
    \right)
    \boldsymbol{\beta}
    \\
    -
    \frac{1}{h}
    \left(
    \boldsymbol{\widetilde{J}}
    \boldsymbol{Q}^{-1}
    \boldsymbol{b}
    +
    \boldsymbol{\widetilde{\phi}}
    \right)^{\top}
    \boldsymbol{\beta}
    -
    \frac{1}{2}
    \boldsymbol{b}^{\top}
    \boldsymbol{Q}^{-1}
    \boldsymbol{b}.
\end{multline}
in which $\boldsymbol{R}$ is a small regularization term commonly used to enhance numerical stability, and
\begin{equation}
\boldsymbol{\widetilde{J}}
=
\begin{bmatrix}
    \boldsymbol{J}_1^n-\mu_1\boldsymbol{J}_{1,1}^t\\
    \vdots\\
    \boldsymbol{J}_1^n-\mu_1\boldsymbol{J}_{1,n_t}^t\\
    \vdots\\
    \boldsymbol{J}_{n_c}^n-\mu_{n_c}\boldsymbol{J}_{n_c,1}^t\\
    \vdots\\
    \boldsymbol{J}_{n_c}^n-\mu_{n_c}\boldsymbol{J}_{n_c,n_t}^t
\end{bmatrix},
\boldsymbol{\widetilde{\phi}}
=
\begin{bmatrix}
    \phi_1\\
    \vdots\\
    \phi_1\\
    \vdots\\
    \phi_{n_c}\\
    \vdots\\
    \phi_{n_c}
\end{bmatrix},
\boldsymbol{\beta}
=
\begin{bmatrix}
    \beta_{1,1}\\
    \vdots\\
    \beta_{1,n_d}\\
    \vdots\\
    \beta_{n_c,1}\\
    \vdots\\
    \beta_{n_c,n_d}
\end{bmatrix}.
\end{equation}

Here, $\boldsymbol{\beta}$ is not only a dual variable but also has a physical interpretation as the contact impulse.
Specifically, the solution of~\eqref{equ.dual_prob_relax} satisfies the following complementarity constraints:
\begin{equation}
\label{equ.dual_prob_relax_sol}
    \boldsymbol{0}
    \leq
    \boldsymbol{\beta}
    \perp
    \frac{1}{h}
    \left(
    \boldsymbol{\widetilde{J}}
    \boldsymbol{Q}^{-1}
    \boldsymbol{\widetilde{J}}^{\top}
    +
    \boldsymbol{R}
    \right)
    \boldsymbol{\beta}
    +
    \left(
    \boldsymbol{\widetilde{J}}
    \boldsymbol{Q}^{-1}
    \boldsymbol{b}
    +
    \boldsymbol{\widetilde{\phi}}
    \right)
    \geq
    \boldsymbol{0}.
\end{equation}

Assuming there exists a positive-definite diagonal matrix that satisfies
$
\boldsymbol{K}(\boldsymbol{q})
=
(
\boldsymbol{\widetilde{J}}
\boldsymbol{Q}^{-1}
\boldsymbol{\widetilde{J}}^{\top}
+
\boldsymbol{R}
)^{-1},
$
we get the closed-form solution of~\eqref{equ.dual_prob_relax} as:
\begin{equation}
    \boldsymbol{\beta}^{+}
    =
    \max
    \left(
    -h
    \boldsymbol{K}(\boldsymbol{q})
    \left(
    \boldsymbol{\widetilde{J}}
    \boldsymbol{Q}^{-1}
    \boldsymbol{b}
    +
    \boldsymbol{\widetilde{\phi}}
    \right),
    \boldsymbol{0}
    \right).
    \label{eq:cf_impulse}
\end{equation}

\textcolor{blue}{The complementarity-free model is then written as:}
\begin{equation}
\label{equ:cf_model}
f_{\mathrm{cf}}:
\begin{cases}
    \boldsymbol{v}^{+}
    =
    \dfrac{1}{h}
    \boldsymbol{Q}^{-1}
    \boldsymbol{b}
    +
    \dfrac{1}{h}
    \boldsymbol{Q}^{-1}
    \boldsymbol{\widetilde{J}}^{\top}
    \boldsymbol{\beta}^{+},
    \\[2mm]
    \boldsymbol{x}_{k+1}
    =
    \boldsymbol{x}_{k}
    +
    h
    \left(
    \boldsymbol{v}_{k}
    +
    \boldsymbol{v}^{+}
    \right).
\end{cases}
\end{equation}
where $\boldsymbol{\beta}$ denotes the contact impulses of all contact points and $\boldsymbol{v}^{+}$ combines the object velocity $\boldsymbol{v}_o$ and robot velocity $\boldsymbol{v}_r$.

\textcolor{blue}{
As discussed in~\cite{jin2024complementarity}, this assumption provides a computationally efficient approximation of the original contact model and enables the closed-form update in~\eqref{eq:cf_impulse}.
}